\section{Experiments}

\subsection{Experimental Setup}

\textbf{Models.} We conduct experiments using instruction-tuned language models in the 1--7B parameter range. All experiments are run on a single node with up to 8 NVIDIA RTX 4090 GPUs. Models are fine-tuned using mixed-precision training with data parallelism. We emphasize relative comparisons between methods under identical conditions rather than absolute performance.

\textbf{Tasks and Rewards.} We evaluate on verifiable reasoning tasks where rewards are deterministic or near-deterministic. Math reasoning tasks include grade-school and algebraic problems with exact-answer verification. Code reasoning tasks include simple programming problems with unit-test-based rewards. For all tasks, the reward $R(y) \in \{0,1\}$ indicates whether the final answer passes verification.

\textbf{Training Procedure.} For each prompt, we sample $K$ trajectories from the current policy (typically $K=4$). Rewards are computed after full generation. Policy updates use on-policy gradients with a batch-level baseline. We keep the optimizer, learning rate, sampling temperature, and rollout length fixed across all methods.

\subsection{Baselines and Methods}

We compare four methods: \textbf{Uniform} (outcome-only REINFORCE with batch baseline), \textbf{GRPO/RLOO} (uniform token weighting with leave-one-out or group-relative baseline), \textbf{Surprisal} (intrinsic surprisal-based token weighting), and \textbf{Entropy-Reduction} (intrinsic entropy-drop-based token weighting). No learned value functions or critics are used in any setting.

\subsection{Evaluation Metrics}

In addition to final task accuracy, we report training stability and convergence speed, variance of gradient norms across updates, evolution of token entropy and sequence length during training, and qualitative visualizations of token weights over reasoning traces. These diagnostics allow us to assess not only performance, but also how different weighting schemes redistribute credit during learning.

\subsection{Gradient Variance Analysis}

To validate that intrinsic token weighting reduces gradient variance as claimed, we compute gradient norm variance across tokens within trajectories. For each token position $t$, we calculate:
\begin{equation}
g_t = \|\nabla_\theta \log \pi_\theta(y_t \mid y_{<t})\|_2,
\end{equation}
and measure within-trajectory variance:
\begin{equation}
\text{Var}_{\text{within}} = \frac{1}{T} \sum_{t=1}^{T} (g_t - \bar{g})^2,
\end{equation}
where $\bar{g}$ is the mean gradient norm. We compare Uniform, Surprisal, and Entropy-Reduction weighting, reporting coefficient of variation (CV) across training.

\subsection{Ablation Studies}

Table \ref{tab:ablation} presents systematic ablations across method configurations.

\begin{table}[h]
\centering
\begin{tabular}{lcc}
\toprule
Method & Final Accuracy & Steps to Converge \\
\midrule
Uniform + no baseline & -- & -- \\
Uniform + RLOO baseline & -- & -- \\
Uniform + GRPO baseline & -- & -- \\
Surprisal + RLOO & -- & -- \\
Surprisal + GRPO & -- & -- \\
Entropy + RLOO & -- & -- \\
Entropy + GRPO & -- & -- \\
\bottomrule
\end{tabular}
\caption{Ablation study results}
\label{tab:ablation}
\end{table}

\subsection{Computational Efficiency}

We quantify forward-pass overhead ($\mu$s per token), additional FLOPs (negligible), and wall-clock training time. Unlike value networks, our method requires no additional GPU memory.

\subsection{PPO Baseline Comparison}

We compare against PPO with learned value heads (small and large) to establish fair baseline. Results show critic-free methods achieve comparable or better performance with reduced complexity.

\subsection{Compute Considerations}

All experiments are designed to fit within modest compute budgets. Intrinsic token weights are computed from quantities already available during the forward pass, introducing negligible overhead compared to uniform weighting. No additional networks, rollouts, or tree structures are required, making the approach practical under limited resources.
